% Student-facing draft, also the standard for model solutions.
% Kept separate for possible inclusion in future homework instructions.
\paragraph{Writing mathematical solutions.}
Write so that a reader can verify each step without having to reconstruct
missing calculations. Prefer formulae to extended prose: organize proofs
around explicit equalities, inequalities, and logical implications, with
short explanations of why the steps are valid.
\begin{enumerate}[leftmargin=*]
\item \textbf{Answer and justify.} Address every assigned part, including
yes/no questions. State the answer and support it with a calculation, proof,
or counterexample, as appropriate.
\item \textbf{Supply the reason at the point of use.} When a step needs
justification, give it: a short clause or an intermediate equality is often
enough. Write out the calculation instead of narrating its steps in words.
Show the identities and matrix operations being used. For example,
if a rotated vector has norm one, explain that the original vector is unit
length and that the rotation preserves norm.
\item \textbf{Connect to the lecture.} Identify the definition, equation,
theorem, or earlier exercise you use. State how its quantities correspond
to yours and check its assumptions. You may use established results without
reproving them unless the exercise asks for their proof.
\item \textbf{Make notation and randomness explicit.} Define new symbols
and say which quantities are fixed and which are random. Define all random
quantities jointly on a probability space, or specify their joint law and
independence assumptions, before using $\P$. Each probability expression must
then have an unambiguous meaning from the notation itself; do not append a
sentence saying which randomness the probability is ``over.'' Use
$\mathbb E[\cdot]$ and explain in words what is being averaged. Identify
conditioning events and justify any independence used; conditioning does
not automatically preserve independence. Omit zero-probability conditioning
cases from averages or explain why their assigned values do not matter.
\item \textbf{Define mathematical objects explicitly.} Specify sets,
events, conditions, and index ranges before using descriptive labels.
Avoid prose placeholders such as ``some index violates the bound.''
For example, define $E_h=\{Z_h>t\}$ for $h\in[H]$ before writing
$\mathbb P(\bigcup_{h=1}^H E_h)\le\sum_{h=1}^H\mathbb P(E_h)$.
Display the set inclusions or other implications used in the argument.
\item \textbf{Check the scope of the argument.} Use the stated assumptions
and cover the allowed boundary and equality cases. Distinguish a claim for
one fixed input from a simultaneous or uniform guarantee. For an asymptotic
claim, state what grows and what is held fixed.
\item \textbf{Prove the kind of claim requested.} An upper bound on error gives a
sufficient guarantee; it does not by itself prove necessity. To establish
an optimum, prove a bound for every allowed choice and give a choice that
attains it. Verify that a construction or counterexample satisfies the
assumptions, and respect which information each decision may use.
\item \textbf{Explain limiting operations.} When exchanging sums,
expectations, integrals, or limits, give the applicable reason, such as a
finite sum, nonnegative summands, or a named theorem whose hypotheses hold.
\item \textbf{End with the requested conclusion.} Relate the calculation
back to the question and, when useful, state its conceptual implication in
one sentence. Distinguish exact values, bounds, and approximations. If an
argument is incomplete, identify the gap explicitly.
\end{enumerate}
